% TikZ 图：开普勒椭圆轨道定律与面积定律
% Day 12：万有引力与天体运动
\begin{tikzpicture}[scale=0.76,>=Latex,line cap=round,line join=round,
  every node/.style={font=\small}]
  % 第一定律：太阳位于焦点
  \begin{scope}[xshift=-3.55cm]
    \coordinate (O) at (0,0);
    \coordinate (S) at (-1.56,0);
    \draw[thick] (O) ellipse (2.0 and 1.25);
    \filldraw[yellow!65!orange,draw=orange!80!black] (S) circle (0.18);
    \node[below=4pt] at (S) {太阳（焦点）};
    \filldraw[blue!55,draw=blue!80!black] (1.35,0.92) circle (0.10);
    \node[above right=2pt] at (1.35,0.92) {行星};
    \filldraw[black] (O) circle (1.2pt);
    \draw[<->,thin] (-2.0,-1.55) -- (2.0,-1.55);
    \node[below=2pt] at (0,-1.55) {$2a$};
    \node[font=\normalsize] at (0,1.80) {椭圆轨道定律};
  \end{scope}

  % 第二定律：相等时间扫过相等面积
  \begin{scope}[xshift=3.45cm]
    \coordinate (O) at (0,0);
    \coordinate (S) at (-1.56,0);
    \draw[thick] (O) ellipse (2.0 and 1.25);
    \filldraw[yellow!65!orange,draw=orange!80!black] (S) circle (0.18);
    \node[below=3pt] at (S) {太阳};

    % 近、远日点附近的两段扫掠区域（示意）
    \fill[red!22]
      (S) -- (-1.86,0.46)
      arc[start angle=158,end angle=202,x radius=2.0,y radius=1.25]
      -- cycle;
    \draw[red,thick] (S) -- (-1.86,0.46);
    \draw[red,thick] (S) -- (-1.86,-0.46);

    \fill[blue!20]
      (S) -- (1.997,-0.065)
      arc[start angle=-3,end angle=3,x radius=2.0,y radius=1.25]
      -- cycle;
    \draw[blue,thick] (S) -- (1.997,-0.065);
    \draw[blue,thick] (S) -- (1.997,0.065);

    \node[red] at (-1.40,0.76) {$\Delta A_1$};
    \node[blue] at (0.75,0.28) {$\Delta A_2$};
    \node[align=center] at (0,-1.75)
      {$\Delta t_1=\Delta t_2$ 时，$\Delta A_1=\Delta A_2$};
    \node[font=\normalsize] at (0,1.80) {面积定律};
  \end{scope}
\end{tikzpicture}
